% ==========================================
\section{Introduzione}
% ==========================================
\subsection{Obiettivi del sistema}
``DeQ: Divide et Quantifica'' è un'applicazione per l'analisi visiva dei mercati e il trading simulato, pensata per analisti, trader e studenti che vogliono libertà nel comporre i propri strumenti.

L'idea di base è superare le dashboard a griglia fissa: al loro posto c'è uno spazio di lavoro infinito, dove si esplorano i dati in tempo reale, si costruiscono i propri schemi collegando dei nodi e si compra e vende in modo simulato.

In particolare, il sistema punta a:

\begin{itemize}[itemsep=-8pt]
    \item \textbf{Analisi modulare:} un catalogo di widget (azioni, grafici, aggregatori) da trascinare e collegare per far scorrere e calcolare i dati in tempo reale.
    \item \textbf{Trading simulato realistico:} un Terminale Broker collegato a dati esterni, per provare acquisti e vendite (con commissioni) senza rischiare denaro vero.
    \item \textbf{Classifica e condivisione:} una Leaderboard dove gli utenti pubblicano le proprie strategie, si confrontano sul guadagno o la perdita delle posizioni e si scambiano recensioni.
    \item \textbf{Salvataggio e prestazioni:} salvataggio automatico della canvas sul database e una cache delle quotazioni per ridurre le chiamate esterne.
\end{itemize}


\subsection{Definizioni}

\arrayrulecolor{SlateGray} % Cambia il colore delle linee della tabella
\begin{tabularx}{\textwidth}{|c|X|}
\hline
\headerTabella{Termine}{Descrizione}
\rigaTabella{Utente (Trader)}{Persona registrata che usa la piattaforma per creare strategie e simulare operazioni.}
\rigaTabella{Workspace (Canvas)}{La tela infinita, navigabile con pan e zoom, dove l'utente dispone e collega i nodi per comporre la sua dashboard.}
\rigaTabella{Widget (Nodo)}{Modulo autonomo (Azione, Grafico, Orologio, Media\ldots) da mettere sulla canvas. Può fare da sorgente o da destinazione dei dati tramite i collegamenti.}
\rigaTabella{Terminale Broker}{La parte per comprare e vendere: listino, carrello, inventario e storico, dove si spende il credito virtuale.}
\rigaTabella{Leaderboard Entry}{La scheda che un utente rende pubblica per la classifica. Il punteggio è il guadagno o la perdita delle posizioni, cioè la differenza tra prezzo attuale e prezzo medio d'acquisto.}
\rigaTabella{Provider Dati}{Servizio esterno (es. Yahoo Finance) da cui il backend prende quotazioni e serie storiche per grafici e listino.}
\rigaTabella{Token JWT}{Il token usato per l'accesso. Grazie a un numero di versione garantisce una sola sessione: un nuovo login disconnette gli altri dispositivi.}
\rigaTabella{Modale}{Finestra che appare sopra la schermata e blocca il resto finché non ci si interagisce (es. per confermare un acquisto).}
\end{tabularx}

% ==========================================
\section{Account di Test e Credenziali}
% ==========================================
Di seguito le credenziali degli account di prova già presenti nel database, per valutare il sistema e i diversi permessi. Per avviare il backend è però necessario prima configurare il proprio file \texttt{server/.env} con le chiavi richieste (vedi \textit{Configurazione dell'ambiente} nel capitolo Architettura).

\begin{table}[H]
\centering
\arrayrulecolor{SlateGray}
\renewcommand{\arraystretch}{1.3}
\begin{tabularx}{\textwidth}{|l|c|c|X|}
\hline
\rowcolor[HTML]{3B3766}
\textcolor{white}{\textbf{Email}} & \textcolor{white}{\textbf{Password}} & \textcolor{white}{\textbf{Ruolo}} & \textcolor{white}{\textbf{Scopo del Test}} \\ \hline
\textbf{\texttt{admin@gmail.com}} & \texttt{adminadmin} & Admin & Account principale. Accesso completo al Pannello Admin per la gestione degli utenti e del listino asset. \\ \hline
\textbf{\texttt{adminriserva@gmail.com}} & \texttt{adminriserva} & Utente & Account standard di riserva, utile come profilo secondario di collaudo. \\ \hline
\textbf{\texttt{carola@gmail.com}} & \texttt{carolauser} & Utente & Account standard. Per provare workspace, acquisti nel Broker e pubblicazione in classifica. \\ \hline
\textbf{\texttt{emanuele@gmail.com}} & \texttt{emanueleuser} & Utente & Account standard. Per provare le interazioni tra utenti, come le recensioni sulle schede altrui. \\ \hline
\textbf{\texttt{daniela@gmail.com}} & \texttt{danielauser} & Utente & Account standard. Per riempire la classifica e vederne l'ordinamento con più utenti. \\ \hline
\textbf{\texttt{michele@gmail.com}} & \texttt{micheleuser} & Utente & Account standard. Per provare l'upgrade a Premium e l'acquisto in blocco di un portafoglio dalla classifica. \\ \hline
\textbf{\texttt{giuseppe@gmail.com}} & \texttt{giuseppeuser} & Utente & Account standard. Per provare i collegamenti tra widget e gli aggregatori di calcolo. \\ \hline
\end{tabularx}
\end{table}